\documentclass[11pt,a4paper]{article}
\usepackage{fontspec}
\setmainfont[AutoFakeBold=2,AutoFakeSlant=.2]{SimSun}
\setsansfont{Microsoft YaHei}
\setmonofont{SimSun}
\XeTeXlinebreaklocale "zh"\XeTeXlinebreakskip=0pt plus 1pt
\usepackage[margin=2.25cm,headheight=15pt]{geometry}
\usepackage{amsmath,amssymb,booktabs,array,longtable,graphicx,xcolor,caption,hyperref,fancyhdr}
\definecolor{ink}{HTML}{16354B}
\hypersetup{colorlinks=true,linkcolor=ink,urlcolor=ink,citecolor=ink,
 pdftitle={量价选股与跨资产风险预算：一项可执行的多资产组合策略},
 pdfauthor={},pdfsubject={共享现金账户、成本后绩效与资产收益贡献的探索性研究}}
\setlength{\parindent}{2em}\setlength{\parskip}{.35em}\linespread{1.13}
\setlength{\tabcolsep}{5pt}\renewcommand{\arraystretch}{1.2}
\renewcommand{\tablename}{表}\renewcommand{\figurename}{图}
\renewcommand{\refname}{参考文献}\renewcommand{\contentsname}{目录}
\captionsetup{font=small,labelfont=bf,labelsep=quad}
\pagestyle{fancy}\fancyhf{}
\fancyhead[L]{\small 量价选股与跨资产风险预算}\fancyhead[R]{\small 策略研究报告}
\fancyfoot[C]{\small\thepage}\renewcommand{\headrulewidth}{.3pt}
\emergencystretch=2em
\newcommand{\note}[1]{\par\vspace{.35em}{\footnotesize\noindent 注：#1\par}}

\newcommand{\reportfigure}[3]{\begin{figure}[htbp]\centering\includegraphics[width=\linewidth]{figures/#1.pdf}\caption{#2}\label{#3}\end{figure}}
\begin{document}
\thispagestyle{plain}
\begin{center}
 {\LARGE\bfseries\color{ink}量价选股与跨资产风险预算}\\[.7em]
 {\large 一项可执行的多资产组合策略}\\[.8em]
 {\small 研究报告公开版\quad 2026年11月5日}
\end{center}
\vspace{.4em}
\begin{center}\textbf{摘要}\end{center}
本文研究如何把股票相对选股能力转化为可执行的多资产组合。策略由两层构成：股票层以滚动训练的LightGBM模型，在历史沪深300成分内进行量价排序；资产层在主动股票、境内标普500 ETF、黄金ETF及国债ETF之间分配风险，每20个交易日依据收缩协方差更新目标权重。所有模块共用现金账户，计入整手、交易费用、滑点、分红与交易限制，避免将独立净值的简单加权误认为可实现收益。

2024—2025年，最终候选RISK-B1成本后累计收益27.80\%、零现金基准夏普2.858、最大回撤3.38\%；同对象、同再平衡频率的固定25\%配置分别为46.13\%、2.642及5.55\%。风险预算提高了该段历史的收益波动比，并以较低收益换取较浅回撤。资产贡献分解显示，黄金与标普500合计贡献19.55个百分点，占主组合净利润约70.32\%；该结果更支持跨资产收益来源与风险管理的解释，而非新增选股alpha。

2026年1月5日至9月21日的175个交易日，以独立100万元现金重新起步，RISK-B1收益2.96\%、夏普0.446、最大回撤8.76\%，固定配置为0.52\%、0.122及10.22\%。相对表现仍较好，但绝对绩效明显减弱。本文以最终策略为主线，归纳财务、趋势、区制、序列、新闻NLP和注意力等探索的取舍，并明确区分训练复现、冻结预测回放与未完成证据。所有评价期均已不同程度用于研究，结论属于探索性历史证据。

\noindent\textbf{关键词：}量价选股；多资产配置；风险预算；共享现金账户；收益归因

\vspace{.5em}
\noindent\textbf{核心判断。}本项目的主要成果是把相对选股、跨资产配置和真实账户约束连接起来，并提供可追溯的结果。RISK-B1是夏普优先目标下的最终候选；固定配置是强对照，绝非已被证明劣于动态模型。本文不以算法数量或复杂度作为策略有效性的依据。

\clearpage
\section{投资问题与策略定位}
\subsection{从相对选股到资本配置}
即使模型总能选出截面上相对较好的股票，也不意味着股票市场在任何阶段都值得充分暴露。初期E00股票策略在2024年取得39.49\%收益，2025年降至8.05\%，提示单一股票风险来源与量价关系存在年度依赖。仅更换预测器，可能仍在承担相同的市场与风格风险。

本策略将决策拆成两个经济问题：第一，在可交易股票中买谁；第二，在股票、海外权益、黄金和利率资产之间承担多少风险。前者由量价排序完成，后者由风险预算完成。风险预算不预测下一段一定上涨的资产，而是根据近期波动及相关性调整资本，使各模块的模型风险份额接近预设目标。该思路与风险贡献分解有关，等风险贡献只是其特殊情形\cite{erc}。

\begin{table}[htbp]\centering\small
\caption{最终策略的两层职责与投资对象}\label{tab:framework}
\begin{tabular}{@{}p{2.3cm}p{5.0cm}p{7.0cm}@{}}\toprule
层次 & 决策 & 规则与职责 \\\midrule
股票选择 & 历史沪深300成分内相对排序 & E00月度滚动训练；目标80只股票，每5交易日检查，前120名保留缓冲 \\
资产配置 & 四个模块的资本目标 & RISK-B1每20交易日更新；股票、标普500、黄金、国债风险预算为45/30/20/5 \\
统一执行 & 共享现金与实际持仓 & 信号收盘形成、次日开盘执行；交易约束可能使实际权重偏离目标 \\
\bottomrule\end{tabular}\end{table}

\subsection{为什么选择这四类风险来源}
主动股票模块保留量价选股研究基础；513500提供境内交易的美国大盘权益敞口；518800提供黄金价格敞口；511010提供人民币国债利率敞口。境内QDII的人民币成交收益同时受海外市场、汇率及溢折价影响，不能等同于美元标普指数收益。国债ETF也有久期与价格风险，不是现金替代品。跨资产相关性并非恒定，更不保证危机期间必然分散。

现金一直允许持有，按零计息处理。最终优化器约束四模块目标资本之和为1，没有单独训练的现金择时器；实际现金主要由整手、费用、未成交及执行顺序形成。因此本策略实现了动态风险分配，尚未实现对所有资产独立判断“是否值得持有”的绝对收益专家系统。

\subsection{评价目标与对照选择}
后期项目明确以成本后夏普为首要指标，同时披露累计收益与回撤。早期股票升级曾要求兼顾弱年收益、回撤和成本，两个阶段的门槛并不完全相同。固定25\%对照在四模块之间每20交易日恢复目标比例，股票内部仍按E00更新；它是持续选股与再平衡的策略，而非买入后完全不动。比较它的目的，是检验额外动态配置层是否值得保留。

\clearpage
\section{数据、信息时间与评价设计}
\begin{table}[htbp]\centering\small
\caption{数据用途与主要来源}\label{tab:data}
\begin{tabular}{@{}p{3.0cm}p{3.1cm}p{8.2cm}@{}}\toprule
数据 & 获取来源 & 用途及边界 \\\midrule
股票行情、历史成分 & 掘金 & 复权量价用于信号，原价用于账户；历史成分按当期使用，129个月并集637只不代表每日637只 \\
ETF风险代理 & Tushare fund\_daily & 构造配置风险输入；成交额原千元统一为人民币元 \\
ETF执行行情 & 东方财富push2his & 境内ETF开收盘成交估值；与风险代理分别保留来源和版本 \\
权益事件 & 供应商、发行人公告等 & 分红、送转、税负与份额独立记账；2026国债9月事件初级来源仍不完整 \\
新闻及业绩预告 & iFinD、比赛邮件 & 仅用于独立探索及覆盖核查；不进入最终E00与B1训练输入 \\
\bottomrule\end{tabular}\end{table}

\subsection{时间切分及研究者可见信息}
\begin{table}[htbp]\centering\small
\caption{各阶段评价口径}\label{tab:periods}
\begin{tabular}{@{}p{3.3cm}p{2.0cm}p{8.9cm}@{}}\toprule
区间 & 账户日数 & 用途及解释 \\\midrule
2020—2022 & 依路线不同 & 配置预测模型等的训练期；E00另按36个月滚动训练，并非统一静态训练 \\
2023 & 242 & 已查看的开发验证，100万元现金起步；用于B1/B2选择 \\
2024—2025 & 485 & 2024年初100万元现金起步，2025继承持仓；多轮查看的探索期 \\
2026-01-05—09-21 & 175 & 每方案各100万元现金重启；已查看背景的事后补充，不是全年或连续跨年 \\
\bottomrule\end{tabular}\end{table}

训练与评价不存在“一律用2020—2024训练”的统一设定。E00每月使用此前36个月且标签已成熟的样本；完整历史重训入口覆盖2020—2025共72个月。为支持首月窗口，底层训练面板保留2017起历史。2026已有九个月模型及冻结评分，后月可使用此前已成熟的2026标签；交付入口只回放这些评分，未封装2026重新训练。

\subsection{防止时间错配的实现约定}
股票特征截至信号日形成，交易最早在下一交易日开盘。训练筛选同时约束信号时间与标签终点，不能因未来标签尚未成熟而删除当前应预测的证券。财务探索按公告可得时间和陈旧度处理，但历史修订链并不完整。新闻实验另设发布时间后第二交易日生效和延迟对照，不把检索得到文本的日期当作原始发布日期。

以上是实现中的时间约束，不会消除研究者反复观察、修改路线和选择候选带来的偏差。本文不把程序上未使用未来标签等同于研究设计上的独立盲测。

\clearpage
\section{股票层：量价排序模型E00}
\subsection{输入与预测目标}
E00采用LightGBM回归器\cite{lgb}。20项输入覆盖收益与趋势、波动与尾部、成交与流动性以及市场敏感度；它是多维量价模型，不能简单称为只买过去涨幅最高股票的动量策略。

\begin{table}[htbp]\centering\small
\caption{E00特征分组及冻结参数}\label{tab:features}
\begin{tabular}{@{}p{3.0cm}p{11.3cm}@{}}\toprule
组别 & 特征或参数 \\\midrule
收益及趋势（8项） & ret1、ret5、ret20、ret60、ret120、momentum20\_skip5、ma20\_distance、ma60\_distance \\
波动及价格路径（5项） & vol20、vol60、downside20、range20、max\_return20 \\
成交及流动性（3项） & volume\_ratio5\_20、log\_amount20、illiquidity20 \\
市场与日内结构（4项） & beta60、residual\_vol60、overnight20、intraday20 \\\midrule
训练窗口与频率 & 滚动36个月，每月一次；各训练日期总权重相同 \\
模型容量 & 15叶、100轮、学习率0.05、最小叶内样本100、L2正则5 \\
可重复设置 & 种子20261006，4线程，deterministic与force\_col\_wise启用 \\
\bottomrule\end{tabular}\end{table}

设$\widetilde O_{i,t}$为复权开盘价，E00的五交易日预测标签及截面目标为
\begin{equation}
 y_{i,t}=\frac{\widetilde O_{i,t+6}}{\widetilde O_{i,t+1}}-1,\qquad
 z_{i,t}=\operatorname{RankPct}_{i}(y_{i,t})-\tfrac12.
\end{equation}
月初拟合仅纳入标签终点不晚于上一个全局可用交易日的样本。排名目标强调相对顺序，既不直接估计市场绝对回报，也不提供股票模块应否满仓的结论。

\subsection{从分数到股票持仓}
账户在当期历史沪深300成分内按分数排序，目标持有80只，前120名作为保留缓冲，每5交易日检查。每次股票预算由外层配置提供，单股买入上限3\%，并服从整手、现金、T+1及停牌、涨跌停等限制。即使目标股票权重相同，高价股票的整手门槛和最低佣金也会影响最终组合。

股票层和资产层采用不同频率：月度更新模型，五交易日检查股票，二十交易日更新资产目标。三种时钟分别对应模型估计、选股执行和风险预算，预测五日收益不意味着每天必须清仓再买。

\subsection{保留E00的理由与证据限度}
E00提供可训练、可排序且已接入账户的股票基线。它在2024年的较好表现和2025年的减弱同时保留。尚缺将E00替换为可交易宽基ETF、再匹配风险与现金占比的完整消融，因此不能单凭最终多资产组合夏普较高，就证明主动选股必不可少或具有稳定独立alpha。

\clearpage
\section{资产层：风险预算与账户执行}
\subsection{风险预算求解}
在每20交易日节点，使用截至当日收盘的最近60个交易日四资产日收益估计协方差，向对角阵收缩50\%，并加$10^{-12}I$作为数值稳定项。历史不足40条或存在非有限值时沿用上期目标；初始目标为四项各25\%。股票风险代理为中证500，实际股票模块为E00，两者存在明确的代理错配。

令$\widehat\Sigma_t$为上述协方差，$w$为资本权重，组合波动及边际分解为
\begin{equation}
 \sigma_p(w)=\sqrt{w^\top\widehat\Sigma_t w},\qquad
 RC_i(w)=\frac{w_i(\widehat\Sigma_t w)_i}{\sigma_p(w)}.
\end{equation}
最终风险预算$b=(0.45,0.30,0.20,0.05)$，按以下目标近似求解：
\begin{equation}\label{eq:opt}
 \min_w\sum_{i=1}^{4}\left(\frac{RC_i(w)}{\sum_j RC_j(w)}-b_i\right)^2
 +0.02\left(\sum_i|w_i-w^{\mathrm{prev}}_i|\right)^2,
\end{equation}
\begin{equation}
 w_i\geq0,\qquad \sum_iw_i=1,\qquad
 w\leq(0.8,0.7,0.6,0.6).
\end{equation}
其中$w^{\mathrm{prev}}$是上一次目标权重，不是当日实际持仓；惩罚项控制目标变化，并非逐笔交易成本的精确模型。SLSQP最多迭代100次，容差$10^{-11}$，求解失败沿用上期目标。边界与惩罚使风险份额不一定精确等于预算。

\subsection{为什么国债资本占比可能很高}
风险贡献与资本占比不同：低波动资产即使只分配少量风险预算，也可能需要较多本金。B1在2024—2025的平均目标资本权重为股票14.16\%、标普14.87\%、黄金14.54\%、国债56.42\%。2026则变为27.46\%、31.99\%、15.46\%、25.09\%。这反映输入风险与约束下的权重变化，不是固定熊市强制买债，也不是准确识别宏观周期的证明。

\subsection{共享现金、费用与公司行为}
每期收盘形成目标，次日开盘处理股票与ETF交易；股票和ETF不能重复使用同一笔现金。标准股票佣金假设0.03\%、最低5元，过户费及印花税按冻结适配器处理，单边滑点5基点。压力情景将费用及滑点加倍，股息税另按持有期记账；ETF按相应产品规则执行。滑点已进入成交价，不从净值再次扣除。

账户净资产定义为$V_t=C_t+\sum_iq_{i,t}P_{i,t}+D_t-T_t$，其中现金、持仓原价市值、股息应收及估计股息税负分别记录。复权仅用于信号，避免再从账户重复计入分红。日线级交易限制无法还原真实盘口排队和全部部分成交。

\clearpage
\section{主结果：收益与风险的明确取舍}
\subsection{2024—2025年共享账户}
表\ref{tab:main}在相同对象和账户框架下比较配置层。RISK-B1标准夏普比固定配置高0.216，回撤浅2.17个百分点，但累计收益低18.33个百分点。因而“最终候选”来自后期夏普优先目标，不是所有指标全面最优。

\begin{table}[htbp]\centering\small
\caption{2024—2025年共享账户绩效}\label{tab:main}
\begin{tabular}{@{}llrrrr@{}}\toprule
方案 & 成本 & 累计收益(\%) & 年化收益(\%) & 夏普 & 最大回撤(\%) \\\midrule
RISK-B1 & 标准 & 27.80 & 13.59 & 2.858 & $-3.38$ \\
固定25\% & 标准 & 46.13 & 21.78 & 2.642 & $-5.55$ \\
趋势倾斜 & 标准 & 40.95 & 19.52 & 2.393 & $-4.96$ \\
Ridge配置 & 标准 & 39.55 & 18.91 & 1.788 & $-8.33$ \\
MLP配置 & 标准 & 39.64 & 18.94 & 1.811 & $-9.42$ \\\midrule
RISK-B1 & 压力 & 27.03 & 13.24 & 2.742 & $-3.44$ \\
固定25\% & 压力 & 44.40 & 21.03 & 2.556 & $-5.60$ \\
趋势倾斜 & 压力 & 38.03 & 18.23 & 2.254 & $-5.05$ \\
\bottomrule\end{tabular}
\note{2024-01-02至2025-12-31，485日，各100万元现金起步，2025继承持仓。零现金基准夏普采用252日约定。Ridge/MLP预测20日原价开盘收益，训练2020—2022、开发验证2023；它们不是与风险预算只差一个算法的严格配对实验。}
\end{table}

\reportfigure{historical_nav_drawdown}{2024—2025年标准成本净值与回撤}{fig:history}
\note{净值以初始100万元归一为1；回撤锚点包含初始资金。曲线由既有共享账户结果重绘，本次修订没有重训策略。}

\clearpage
\subsection{开发验证、年度表现与执行暴露}
2023年的已查看开发验证中，B1夏普2.837高于B2的2.776，因此在两个风险候选之间保留B1；固定25\%为2.662。它是一年开发比较，不能作为独立确认性证据。历史年度收益与执行记录见表\ref{tab:annual}和表\ref{tab:execution}。

\begin{table}[htbp]\centering\small
\caption{主方案分期收益与2023开发验证}\label{tab:annual}
\begin{tabular}{@{}llrrl@{}}\toprule
区间 & 指标 & RISK-B1 & 固定25\% & 说明 \\\midrule
2023 & 区间收益(\%) & 8.49 & 12.52 & 独立现金起点 \\
2023 & 零现金夏普 & 2.837 & 2.662 & 已用于开发 \\
2023 & 最大回撤(\%) & $-1.74$ & $-2.50$ & 242日 \\\midrule
2024 & 年度收益(\%) & 17.55 & 24.64 & 连续账户首年 \\
2025 & 年度收益(\%) & 8.72 & 17.23 & 继承2024持仓 \\
\bottomrule\end{tabular}\end{table}

\begin{table}[htbp]\centering\small
\caption{2024—2025年标准账户执行与平均实际暴露}\label{tab:execution}
\begin{tabular}{@{}lrr@{}}\toprule
指标 & RISK-B1 & 固定25\% \\\midrule
成交笔数 & 2,082 & 2,422 \\
累计成交额（万元） & 530.38 & 749.65 \\
显式交易费用（元） & 11,667.54 & 13,906.71 \\
已付股息税（元） & 788.49 & 1,729.18 \\
平均实际股票占比(\%) & 8.08 & 17.99 \\
平均实际ETF占比(\%) & 81.14 & 71.89 \\
平均实际现金占比(\%) & 10.78 & 10.13 \\
\bottomrule\end{tabular}
\note{实际暴露为每日市值/NAV的平均，不能用调仓时目标权重均值替代。滑点已含成交价，费用与股息税分别列示。累计成交额不是年化换手率。}
\end{table}

平均目标股票权重14.16\%最终只形成8.08\%的实际股票暴露，说明配置求解不是账户的全部。股票整手、持仓缓冲及可交易限制会造成预算未充分使用。仅将四条收益序列乘以目标权重，无法复刻本策略现金与费用路径。低实际股票暴露也限制了把高夏普解释为选股能力的空间。

固定25\%没有被假定为最优。曾在2020—2022代理路径上评价132组预设且去重的固定比例，按年度夏普均值减0.5倍年度夏普标准差评分，并过滤非正年度收益。两组候选在2023真实账户夏普2.292、2.267，未胜固定配置2.662。这个有限搜索说明简单对照有竞争力，不证明不存在更优比例。

\clearpage
\section{2026补充检验：相对改善与绝对减弱}
2026账户于1月5日各以100万元现金起步，截至9月21日共175日，不承接2025年末持仓。它不是2024—2026连续净值，也不是2026全年、实盘或独立盲测。策略参数与配置公式沿用冻结版本。

\begin{table}[htbp]\centering\small
\caption{2026年175日现金重启绩效}\label{tab:2026}
\begin{tabular}{@{}llrrr@{}}\toprule
方案 & 成本 & 区间收益(\%) & 零现金夏普 & 最大回撤(\%) \\\midrule
RISK-B1 & 标准 & 2.96 & 0.446 & $-8.76$ \\
固定25\% & 标准 & 0.52 & 0.122 & $-10.22$ \\
RISK-B1 & 压力 & 2.78 & 0.433 & $-8.72$ \\
固定25\% & 压力 & 0.26 & 0.091 & $-10.18$ \\
\bottomrule\end{tabular}
\note{累计收益没有年化为“全年收益”；夏普采用175个日收益和252日惯例，仅为短样本描述。标准与压力各自重新执行账户，不能将二者差解释为固定持仓路径下的纯费用扣除。}
\end{table}

\reportfigure{extension_nav_drawdown}{2026年现金重启的标准成本净值与回撤}{fig:2026}

相对固定配置，B1收益高2.44个百分点、夏普高约0.324、回撤浅约1.46个百分点；但相较自身历史夏普2.858、回撤3.38\%，其绝对表现显著减弱。不能因对照更弱而忽略策略本身的回撤与收益停滞，也不能仅凭此将下降归因于某个新闻事件或宏观因素。

\clearpage
\subsection{实际仓位与压力情景的解释边界}
图\ref{fig:weights}展示标准账户每日实际市值，而非优化器给出的目标。2026 B1平均实际股票15.78\%、标普32.19\%、黄金15.26\%、国债20.94\%、现金15.82\%。该配置与2024—2025明显不同，风险预算具有动态调整，但不会自动及时避开所有下跌。

\reportfigure{actual_weights_2026}{2026年RISK-B1标准账户实际资本分布}{fig:weights}
\note{图中不强制归一化，市值及现金占比之和与1的微小偏差来自股息应收和估计税负，观测范围约99.865\%—100.015\%；不将这些调整伪装成可用现金。}

\begin{table}[htbp]\centering\small
\caption{2026标准与压力账户的执行差异}\label{tab:stress}
\begin{tabular}{@{}llrrrr@{}}\toprule
方案 & 成本 & 成交笔数 & 费用（元） & 平均现金(\%) & 平均股票(\%) \\\midrule
RISK-B1 & 标准 & 1,064 & 6,281.47 & 15.82 & 15.78 \\
固定25\% & 标准 & 963 & 5,571.59 & 15.32 & 13.62 \\
RISK-B1 & 压力 & 1,060 & 12,407.74 & 39.91 & 15.85 \\
固定25\% & 压力 & 957 & 10,955.35 & 39.27 & 13.65 \\
\bottomrule\end{tabular}\end{table}

压力情景的平均现金从约16\%升至约40\%，是需要单独说明的重要现象。账本恒等式与费用合计通过并不解释这一变化的全部机制；已有记录不足以把它完全归结为成本、整手或某一具体执行环节。本报告保留原输出，将其列为进一步审计事项。因此压力回撤略浅不能解释为成本增加有利，标准与压力的接近也不足以证明策略对所有摩擦稳健。

最终候选中尚无独立新闻择时器。2026表现减弱既可能与风险环境变化有关，也可能受收益源、代理错配和账户暴露影响；只有逐项对照才能区分，不能用“宏观未考虑”代替检验。

\clearpage
\section{收益来源：分散贡献与风险管理的边界}
以每个账户初始100万元为分母，将ETF交易、持仓与权益事件重建为资金损益贡献；股票及未分拆项取账户总损益扣除三ETF后的差额。后者含股票、费用、税负、未分配项及潜在上游误差，不能称为选股alpha。现金按原假设零计息。

\begin{table}[htbp]\centering\small
\caption{对初始资金的收益贡献（百分点）}\label{tab:attribution}
\begin{tabular}{@{}lrrrr@{}}\toprule
组件 & B1历史 & 固定历史 & B1补充期 & 固定补充期 \\\midrule
黄金ETF & 12.5212 & 22.2514 & 0.4377 & $-1.2770$ \\
标普500 ETF & 7.0248 & 13.1840 & 2.9937 & 2.1933 \\
国债ETF & 3.4131 & 1.5816 & 0.3657 & 0.3657 \\
股票及未分拆项（差额） & 4.8363 & 9.1081 & $-0.8402$ & $-0.7603$ \\\midrule
合计 & 27.7953 & 46.1251 & 2.9569 & 0.5216 \\
\bottomrule\end{tabular}
\note{历史为2024—2025；补充期为2026年175日现金重启。该表不是各资产自身收益率，亦非Brinson行业归因。三ETF独立重建，股票剩余项的对账不构成独立验证。}
\end{table}

\reportfigure{component_contributions}{两段账户的组件贡献分布}{fig:contribution}
\note{星号表示股票及未分拆项按差额计算。两个面板纵轴尺度不同，以便阅读弱收益期；不能按柱高直接比较两个时期的量级。}

历史B1的黄金与标普共贡献19.55个百分点，占净利润约70.32\%；国债实际平均资本约52.42\%，但只贡献3.41个百分点。资本占比、利润贡献与风险贡献是不同概念。2026 B1相对固定配置多2.4353个百分点，其中黄金损益差贡献1.7148、标普0.8004、国债相同、股票剩余项拖累0.0799。优势集中在黄金和标普的持有路径，尚不能全部归因于风险预算的预测能力。

\clearpage
\section{研究过程的取舍：哪些证据改变了方案}
\subsection{股票失效诊断推动了风险来源扩展}
股票基线2024年与2025年的分化首先引出财务、状态和趋势路线，而非立即接受某种复杂模型。事后暴露回归中，以境内沪深300 ETF为代理，E00的beta由2024年的0.6532变为2025年的0.5442，日收益拟合$R^2$由0.7816变为0.6064；2025截距与零无明显区别。它支持股票关系减弱的描述，不能识别行业、风格与选股技能的因果贡献。2025年第三季度负IC与偏低beta、低波动暴露构成研究线索，不是全年表现的唯一原因。

随后扩大资产风险来源，固定配置成为强对照。研究重心因而从“再增加一个选股因子”转向“同样信号如何分配资本并通过账户落地”。这一调整是最终方案的主要金融逻辑。

\subsection{复杂方法为什么未必带来增量}
新增模块只有在信息、期限、目标及执行上形成有效增量，才可能提高组合质量。输入与原有量价重复时，更大模型可能主要增加估计误差；财务慢变量用于五日标签可能期限不匹配；区制划分可能减少专家样本并造成切换滞后；资金比例预测还要经约束和成交转换，误差会进入实际仓位。上述是与结果一致的机制假说，未经严格消融不能写成已证实根因。

\begin{table}[htbp]\centering\small
\caption{关键取舍与最终方案的关系}\label{tab:lessons}
\begin{tabular}{@{}p{2.5cm}p{5.6cm}p{6.2cm}@{}}\toprule
方向 & 观察事实 & 对最终方案的影响 \\\midrule
财务、期限扩展 & JF20未胜同期限P20；延长期限未消除年度减弱 & 保留量价基线，暂不增加未经证实的信息层 \\
趋势与区制 & T01在2025较好，S01收益改善伴随更深回撤；RG01弱年仍弱 & 互补性值得保留，不直接等同可识别的切换优势 \\
阶段单选 & 主规则未胜简单对照；敏感性局部较好 & 不把事后最佳窗口替代原主规则 \\
新闻与注意力 & NLP当前指数目标未证实增量；注意力局部优于配对MLP & 承认局部证据，不把模型复杂度当作总收益优势 \\
多资产风险配置 & 相同对象下B1历史夏普较高、收益较低 & 作为夏普优先候选，并列强对照与2026退化结果 \\
\bottomrule\end{tabular}\end{table}

\subsection{项目优势与尚未解决的问题}
优势在于明确分层、保留强基准、接入共享现金及权益事件、并把归因与负结果放回同一证据链。尚未解决的是：多次探索后的选择偏差、风险代理与实际股票错配、配置与执行造成的未投资资金，以及未独立验证的收益来源。更高夏普不自动代表更好股票预测，也不代表更适合以累计收益排名为唯一目标的实时比赛。

\clearpage
\section{结论与可复现交付}
本项目形成了“相对选股—风险预算—统一执行”的可复核策略。2024—2025年的结果支持RISK-B1在收益、波动和回撤之间形成不同于固定配置的取舍，收益贡献主要来自跨资产模块。2026补充期仍有相对优势，但低夏普与更深回撤表明该优势不能外推为稳定高收益。现阶段最准确的定位是具有账户和证据基础的探索性多资产研究，而非已确认的普适择时系统。

\subsection{后续最有价值的验证}
优先补充风险匹配的固定配置对照及“E00换成可交易指数”的消融，以判断优势来自选股、资产分配还是风险压低；审计压力情景的现金变化与目标落地；再在未查看时期冻结测试。若继续NLP，应让主题信息对应行业或资产暴露，并保留价格基准、发布时间和延迟对照，而非仅把更多文本加进同一个短期指数标签。以上为后续任务，本次修订没有运行这些新实验。大量尝试后的选择偏差与回测过拟合风险需要单独处理\cite{pbo}；本项目未计算PBO或证明显著性。

\begin{table}[htbp]\centering\small
\caption{交付的三层复现能力}\label{tab:reproduction}
\begin{tabular}{@{}p{2.7cm}p{5.7cm}p{5.9cm}@{}}\toprule
层次 & 可执行内容 & 所需输入与边界 \\\midrule
公开图表再现 & 重画本文四幅结果图、月度收益及指标核对 & 仓库内衍生净值、权重与贡献数据；无需行情账户 \\
冻结预测回放 & 历史六账户和2026四账户，重算配置及交易 & 需自备合法本地输入包；不是只读取汇总，也不是模型重训 \\
历史完整训练 & 2020—2025共72个月E00拟合，再运行历史账户 & 需训练面板与全部账户输入；2026重新训练未封装 \\
\bottomrule\end{tabular}\end{table}

仓库入口为README，完整运行说明为\texttt{代码/README\_使用说明.md}；本文源文件、衍生图表数据、制图脚本和证据清单位于\texttt{docs/strategy-report/}。行情、成分、事件和新闻供应商原始资料不随公开仓库分发；公开代码不能凭空重建未分发输入。主要旧实验的本地恢复范围另见快速复现指南，不能将“有历史摘要”写成“所有分支可从公开仓库一键重训”。

本修订版重组叙述、补图和已有归因，未改参数、训练或原账户结果；原始提交PDF与正式ZIP保留原哈希。方法、数据来源和结果表逐一对应，避免以新的写作日期冒充新的策略验证。

\clearpage
\appendix
\section{主要探索路线与证据状态}
以下记录主要尝试，区分完成账户、仅信号评价、外部返回及未实现。它们的时间、对象和更新频率不完全相同，不将所有夏普放进同一个排名。

{\small
\begin{longtable}{@{}p{2.0cm}p{4.9cm}p{7.3cm}@{}}
\caption{研究路线索引与结论范围}\label{tab:routes}\\\toprule
路线 & 实际实现 & 结果与限制 \\\midrule\endfirsthead
\multicolumn{3}{c}{表\thetable（续）}\\\toprule
路线 & 实际实现 & 结果与限制 \\\midrule\endhead
基线与财务 & 等权B00、规则B01、E00/E10、Qlib适配、JF05/P20/JF20 & 历史账户完成。JF20相对P20两年年化低2.264个百分点，2025低2.622个百分点；未过原升级门槛，不否定财务因子的一般用途。 \\
趋势与条件 & T01中期趋势、S01条件切换、E00/T01双袖 & T01在2025收益17.34\%高于E00的8.05\%，但两年回撤18.64\%。S01两年回撤16.70\%。双袖独立资金记账，不等于两个百万账户收益平均。 \\
风险仓位 & R01动态12\%波动目标及开发期固定仓位 & 动态两年收益32.58\%、回撤8.71\%；固定25.57\%、5.05\%。改变风险暴露有作用，不能直接称增加alpha。 \\
HMM与序列 & RG01概率专家；SQ01冻结ESN状态加Ridge & RG01两年收益55.80\%，但2025仅5.27\%，未解决弱年目标，2026账户未完成。SQ01 2026标准收益$-6.15\%$、夏普$-0.753$，不支持替换。 \\
宏观主题NLP & 六主题文本特征、价格基线及延迟文本 & 2025共2349条检索片段、245个文章日期；1—5月训练、6月隔离、7—12月测试，预测沪深300五日开盘收益。未证实当前目标增量，不代表新闻不能影响油价或行业；2026新闻交易未完成。 \\
时间注意力 & 20日序列、20特征、1层32维2头，与同序列MLP三种子配对 & 注意力标准夏普0.859/1.460/0.968，MLP为0.703/0.655/0.862。平均差约0.356，但块区间跨零；三种子不是独立市场样本。 \\
静态匹配与融合 & N01匹配树、N02股票分数融合 & 静态树标准夏普约0.847；融合三种子约0.884/1.327/0.918，10月14日已用保存工件补预测并回放账户；不等于完整重训。 \\
跨资产注意力 & N03四ETF网络，沪深300替代主动股票 & 固定夏普约2.439、MLP三种子2.913/2.327/2.321、注意力1.219/1.461/1.826；10月14日已由权重及重建标准化回放14账户；原预测数组字节未恢复，训练短且对象不同，不作架构胜负依据。 \\
阶段与现金门 & 股票单选、股票ETF单选、新四ETF价格现金门 & 详见附录B；主规则未稳定胜简单对照。原120/40设想与后续不同评分规则不可混称。 \\
配置搜索 & 132组固定权重、Ridge/MLP预测配置、风险预算、趋势倾斜 & 有限搜索未使固定权重胜验证对照；预测配置未胜同对象固定配置。B1符合后期夏普优先取舍。 \\
其他讨论 & BLO、DRO、HRP、宏观强化学习、直接期货、部分资金资料 & 未完成完整账户检验，不计作失败；未准入的纳指和工业金属路线不得列成已完成策略。 \\
\bottomrule\end{longtable}}

\clearpage
\section{补充账户对照与指标定义}
\begin{table}[htbp]\centering\small
\caption{主要股票路线历史结果（标准成本）}\label{tab:stock}
\begin{tabular}{@{}lrrrr@{}}\toprule
模型 & 2024收益(\%) & 2025收益(\%) & 两年收益(\%) & 最大回撤(\%) \\\midrule
E00量价 & 39.49 & 8.05 & 50.72 & $-8.72$ \\
E10状态特征 & 36.83 & 6.41 & 45.60 & $-9.12$ \\
P20长期标签 & 38.79 & 7.50 & 49.20 & $-8.29$ \\
JF20财务增强 & 37.27 & 4.88 & 43.97 & $-8.82$ \\
T01趋势 & 12.43 & 17.34 & 31.93 & $-18.64$ \\
S01条件 & 28.45 & 13.48 & 45.77 & $-16.70$ \\
RG01隐状态 & 47.99 & 5.27 & 55.80 & $-8.75$ \\
SQ01序列 & 21.17 & 5.96 & 28.39 & $-9.04$ \\
\bottomrule\end{tabular}
\note{2024—2025连续股票现金账户，与正文股票加三ETF账户不同。不同风险暴露下的收益差不是严格选股alpha。2026现金重启E00/T01/S01/E10/P20标准收益分别为$-4.31$/$-4.49$/$-1.86$/$-5.48$/$-3.16$\%。}
\end{table}

\begin{table}[htbp]\centering\small
\caption{阶段选择及价格过滤的局部比较}\label{tab:switch}
\begin{tabular}{@{}p{3.6cm}p{2.3cm}p{2.7cm}p{5.0cm}@{}}\toprule
对象与规则 & 主规则夏普 & 对照夏普 & 证据范围 \\\midrule
股票20日主切换 & 0.531 & S01：0.733 & 2025—2026共418日连续账户；40日敏感性0.752，未跑压力 \\
股票ETF阶段单选 & 0.821 & 同候选固定：1.154 & 同期标准账户；主规则回撤12.78\%，对照8.38\% \\
新四ETF价格现金门 & 1.256 & 同对象固定：1.276 & 60日机会分数、40日节点；回撤较浅，夏普差未过门 \\
\bottomrule\end{tabular}
\note{新四ETF以中证500、标普500、黄金、国债构成，不是正文的E00主动股票。上述连续2026年度段不能与本文175日现金重启结果混用。}
\end{table}

设初始净值$V_0$，$N$个交易日日收益$r_t=V_t/V_{t-1}-1$：
\begin{align}
 R &= V_N/V_0-1,& R_{\rm ann}&=(V_N/V_0)^{252/N}-1,\\
 S_0&=\sqrt{252}\,\overline r/s(r),&
 \mathrm{MDD}&=\min_{0\le t\le N}\left(\frac{V_t}{\max_{0\le s\le t}V_s}-1\right).
\end{align}
$s(r)$为样本标准差。本文夏普参考现金收益0，不是市场超额夏普，也未使用下载的无风险利率；回撤表中以负号列示，文字常用损失幅度。年度收益按跨年净值衔接；收益贡献以原始初始资金为分母，不能与年度收益直接相加。

\clearpage
\section{月度收益与来源索引}
\begin{table}[htbp]\centering\small
\caption{标准成本月度收益（\%）}\label{tab:monthly}
\begin{tabular}{@{}lrrlrrlrr@{}}\toprule
月份 & B1 & 固定 & 月份 & B1 & 固定 & 月份 & B1 & 固定 \\\midrule
2024-01 & 1.22 & 1.33 & 2025-01 & 0.83 & 1.15 & 2026-01 & 2.35 & 3.47 \\
2024-02 & 1.95 & 2.58 & 2025-02 & -1.09 & -1.14 & 2026-02 & -1.19 & -0.78 \\
2024-03 & 3.51 & 3.80 & 2025-03 & 0.36 & 0.77 & 2026-03 & -4.58 & -5.12 \\
2024-04 & 0.08 & 0.17 & 2025-04 & 1.37 & 1.68 & 2026-04 & 4.24 & 2.84 \\
2024-05 & 0.86 & 1.41 & 2025-05 & 0.67 & 1.55 & 2026-05 & 1.06 & 0.35 \\
2024-06 & 1.67 & 2.14 & 2025-06 & 0.83 & 1.45 & 2026-06 & -2.86 & -4.01 \\
2024-07 & 0.59 & 0.94 & 2025-07 & 0.93 & 1.85 & 2026-07 & 0.79 & 0.83 \\
2024-08 & -0.03 & 0.00 & 2025-08 & 0.25 & 0.71 & 2026-08 & 4.27 & 4.19 \\
2024-09 & 2.09 & 3.13 & 2025-09 & 1.99 & 3.55 & 2026-09 & -0.80 & -0.84 \\
2024-10 & 2.72 & 3.69 & 2025-10 & 2.07 & 3.82 & -- & -- & -- \\
2024-11 & 1.14 & 2.69 & 2025-11 & 0.24 & 0.75 & -- & -- & -- \\
2024-12 & 0.55 & 0.45 & 2025-12 & -0.00 & -0.04 & -- & -- & -- \\
\bottomrule\end{tabular}
\note{B1与固定分别指RISK-B1及固定25\%。2024—2025连续；2026重新以现金起步，一月从100万元计算，九月仅截至21日。表由日净值生成；未出现的月份为无评价，不是零收益。}
\end{table}


\noindent\textbf{证据定位。}以下路径相对于仓库根目录：
\begin{itemize}\small
\item 主规则：代码/src/run\_dynamic\_risk\_trend\_v1.py；训练：代码/train\_e00.py。
\item 补充账户：results/最终组合2026/；路线记录：results/历史实验摘要/。
\item 归因说明：docs/收益归因技能旧比赛实测\_20261021.md。
\end{itemize}
本报告目录的\texttt{figure\_sources.json}记录原输入路径与SHA-256，\texttt{figure\_checks.json}记录衍生文件及数值核对。旧版报告保留未在正文展开的实验明细。

\begingroup\small\raggedright
\begin{thebibliography}{9}
\bibitem{lgb} Ke G, Meng Q, Finley T, et al. LightGBM: A Highly Efficient Gradient Boosting Decision Tree[C]. Advances in Neural Information Processing Systems, 2017, 30. \url{https://proceedings.neurips.cc/paper_files/paper/2017/hash/6449f44a102fde848669bdd9eb6b76fa-Abstract.html}.
\bibitem{erc} Maillard S, Roncalli T, Teiletche J. The Properties of Equally Weighted Risk Contribution Portfolios[J]. The Journal of Portfolio Management, 2010, 36(4): 60--70. DOI: 10.3905/jpm.2010.36.4.060. Working paper: \url{https://papers.ssrn.com/sol3/papers.cfm?abstract_id=1271972}.
\bibitem{pbo} Bailey D H, Borwein J M, López de Prado M, Zhu Q J. The Probability of Backtest Overfitting[J]. The Journal of Computational Finance, 2017, 20(4): 39--69. Author manuscript: \url{https://www.davidhbailey.com/dhbpapers/backtest-prob.pdf}.
\end{thebibliography}\endgroup
\end{document}
